\documentclass[12pt]{jlreq}
\usepackage{amsmath, amssymb}

\newcommand{\C}{\mathbb{C}}
\newcommand{\R}{\mathbb{R}}
\newcommand{\Z}{\mathbb{Z}}
\newcommand{\Tr}{\operatorname{Tr}}
\newcommand{\Impart}{\operatorname{Im}}
\newcommand{\AF}{\operatorname{AF}}
\newcommand{\AX}{\operatorname{AX}}
\newcommand{\GL}{\operatorname{GL}}

\begin{document}

測度空間 \((X,\mu)\) について次の命題を考える。

\((*)\) \(X\) 上の実数値可測関数 \(f(x)\) で，ほとんどいたるところ \(0\le f(x)<\infty\) であり，
\[
  \int_X f(x)\,d\mu=\infty
\]
となるものが任意に与えられたときに，次の \(4\) 条件を満たす \(X\) 上の可測関数列
\(\{f_n\}_{n=1,2,3,\ldots}\) が存在する。

\begin{enumerate}
  \item[(a)] すべての \(n\) について，ほとんどいたるところ \(0\le f_n(x)\le f(x)\) である。
  \item[(b)] ほとんどいたるところ \(\lim_{n\to\infty}f_n(x)=0\) となる。
  \item[(c)] \(\displaystyle \lim_{n\to\infty}\int_X f_n(x)\,d\mu=\infty\) である。
  \item[(d)] すべての \(n\) について \(\displaystyle \int_X f_n(x)\,d\mu<\infty\) である。
\end{enumerate}

これについて次の問いに答えよ。

\begin{enumerate}
  \item \((X,\mu)\) が閉区間 \([0,1]\) とその上の Lebesgue 測度であるとき，上の命題 \((*)\) を証明せよ。
  \item \((X,\mu)\) が実数全体の集合 \(\R\) とその上の Lebesgue 測度であるとき，上の命題 \((*)\) を証明せよ。
\end{enumerate}

\end{document}